\documentclass[11pt]{article}
\usepackage[utf8]{inputenc}
\usepackage[T1]{fontenc}
\usepackage[brazilian]{babel}
\usepackage{amsmath, amssymb, amsthm}
\usepackage[margin=2.6cm]{geometry}
\usepackage{booktabs}
\newcommand{\code}[1]{\texttt{\small #1}}
\newcommand{\Z}{\mathbb Z}
\newcommand{\R}{\mathbb R}
\newcommand{\GF}{\mathrm{GF}}
\theoremstyle{plain}
\newtheorem{teorema}{Teorema}
\newtheorem{principio}{Princípio}
\theoremstyle{remark}
\newtheorem{obs}{Observação}

\title{\textbf{O Viveiro}\\[3pt]
\large cada cruzamento dá um pássaro que voa sozinho}
\author{}\date{}

\begin{document}
\maketitle

\begin{abstract}
\noindent Os corpos $\R^n$ do \code{teoria.tex} não são uma lista: são um \textbf{viveiro}. Cada
espécie cruza com outra e o resultado é um terceiro corpo --- que \emph{voa sozinho}, isto é, é ele
próprio um corpo, e não apenas um espaço com a dimensão certa. Esta distinção é a única coisa que este
documento afirma, e é ela que separa a operação certa das duas que parecem servir e não servem: a soma
direta \emph{nunca} voa (tem divisor de zero) e o produto tensorial voa \emph{só} entre espécies
primas entre si. O que voa sempre é o \textbf{cruzamento} $\R^a\vee\R^b=\R^{\mathrm{lcm}(a,b)}$, o
menor corpo em que os dois pais cabem. Ele é comutativo, associativo, \textbf{idempotente} e tem
neutro --- e é a idempotência que faz disto um viveiro e não uma fábrica. Tudo medido em
\code{tools/viveiro.c}, exato, em inteiros, por varredura exaustiva: resíduo $0$.
\end{abstract}

\section{A pergunta que faltava}

O \code{corpodecorpos.c} mediu que a soma direta soma as dimensões e que o produto as multiplica.
Nenhuma das duas medidas respondia ao que importa, e a pergunta certa é de outra ordem:

\begin{center}
\emph{o filhote voa?}
\end{center}

Ou seja: o resultado do cruzamento é ele próprio um \textbf{corpo} --- fechado, com inverso para todo
não-nulo --- ou é apenas um espaço vetorial com a contagem de eixos correta? Ter a dimensão certa é
barato. Voar é que é a condição.

\section{Duas candidatas que não servem}

\begin{teorema}[a soma direta nunca voa]
Em $\R^a\oplus\R^b$, os elementos $(1,0)$ e $(0,1)$ são ambos não-nulos e o seu produto é $(0,0)$.
Logo há divisor de zero, e nenhum dos dois admite inverso. Medido para todo par $a,b\le3$: sem
exceção. \emph{A dimensão fecha em $a+b$; a estrutura não fecha.}
\end{teorema}

\noindent É o retrato exato do que a soma direta faz: põe dois pássaros na mesma gaiola. Não nasce
pássaro novo --- nasce uma gaiola com dois dentro, e cada um continua a ser o que era.

\begin{teorema}[o tensorial voa se e só se as espécies são primas entre si]
$\R^a\otimes\R^b$ tem dimensão $a\cdot b$, e é corpo \emph{se e só se} $\gcd(a,b)=1$. Medido por
força bruta --- procurando divisor de zero na definição, não por critério: para $a=b=2$ há divisor e
não voa; para $(2,3)$, $(3,2)$, e todos os coprimos testados, não há e voa.
\end{teorema}

\begin{obs}[por que o tensorial \emph{parece} funcionar]
Quando $\gcd(a,b)=1$ vale $a\cdot b=\mathrm{lcm}(a,b)$, e aí o tensorial \textbf{coincide} com o
cruzamento. Quem só testar espécies primas entre si concluirá que o produto é a lei do viveiro --- e
estará certo no que viu e errado no geral. Quando as espécies partilham divisor, o tensorial
reparte-se em cópias e o divisor de zero aparece.
\end{obs}

\section{A lei do viveiro}

\begin{teorema}[o cruzamento, e o filhote voa]
Cruzar $\R^a$ com $\R^b$ dá
\[
  \R^a\vee\R^b=\R^{\,\mathrm{lcm}(a,b)} ,
\]
e o filhote \textbf{voa}: é corpo, contém os dois pais, e é o \emph{menor} corpo que os contém.
Medido: para todo par testado, $\R^{\mathrm{lcm}}$ tem inverso para todo não-nulo; a contenção vale
porque $\R^x\subset\R^n\iff x\mid n$; e nenhum $n<\mathrm{lcm}(a,b)$ é divisível por $a$ e por $b$.
\end{teorema}

\noindent Nem a soma nem o produto: a \textbf{junção}. E ela existe para todo par, sem exceção --- não
há espécies incruzáveis no viveiro.

\begin{teorema}[o viveiro é fechado, e é semirretículo]
O cruzamento é comutativo, associativo, \textbf{idempotente} ($\R^a\vee\R^a=\R^a$) e tem neutro
($\R^1=\Z_p$). E é fechado: nos cruzamentos examinados, nenhum filhote deixa de voar e nenhum sai do
viveiro. Por isso se pode cruzar de novo, e de novo: o filhote é da mesma espécie de coisa que os
pais.
\end{teorema}

\begin{obs}[a idempotência é o que faz disto um viveiro e não uma fábrica]
Uma fábrica combina duas peças e produz sempre algo maior. Aqui, cruzar uma espécie \emph{consigo
mesma} devolve ela mesma --- não há descendência de um só. Só o encontro do \textbf{diferente} gera
espécie nova, e o quanto de novo é exatamente $\mathrm{lcm}(a,b)/a$: o que o filhote tem e o pai não
tinha. \emph{A novidade mede-se, e ela vem da diferença.}
\end{obs}

\section{Onde isto se encaixa}

O cruzamento é a junção do retículo de subcorpos medido em \code{tools/corpodecorpos.c} --- ordenado
por divisibilidade, com encontro $\gcd$ e junção $\mathrm{lcm}$. O que este documento acrescenta é que
a \emph{junção} é a operação de criação, e o \emph{encontro} não: $\R^a\cap\R^b=\R^{\gcd}$ é o que os
dois pais já tinham em comum, e não gera nada.

E a dualidade de \code{tools/pontryagin.c} troca as duas: $\mathrm{dual}(\gcd)=\mathrm{lcm}$ dos
duais. Logo, do outro lado do espelho, criar é intersectar. Não são duas leis --- é uma, vista de
cada lado.

\begin{principio}[o critério]
\emph{Ter a dimensão certa não basta; a condição é voar.} Ao propor uma operação entre corpos, a
pergunta não é que tamanho tem o resultado, mas se o resultado é da mesma espécie que as entradas. As
duas operações que falham aqui --- a soma e, fora dos coprimos, o produto --- falham exatamente por
não passarem nesse teste, e as suas dimensões estão corretas nos dois casos.
\end{principio}

\end{document}
